In Tables~\ref{tab:combi} and~\ref{tab:combi:2},  
global Orbiter commands for Combinatorics are summarized.
\orbitertableofcommands{combinatorics_1}
\orbitertableofcommands{combinatorics_2}

\bigskip



The command 

\sourcecode{random_permutation_12}

creates a random permutation in $\Sym(12).$ 
The permutation is written to the csv file named in the command. 
The command seeds the pseudo-random number generator with a number (here, the first few digits of $\pi$), 
just to put the random number generator into a state different from the zero state. 
On some systems, this may prevent that the sequence of random numbers starts with 0. 


\bigskip





The command 

\sourcecode{Sym_10_conj_classes}

produces a list of the conjugacy classes of $\Sym(10).$ 
The list is written to a csv file. 


\bigskip


The next command computes the character table of the symmetric group $\Sym(4)$:

\sourcecode{Char_Sym_4}

The command produces the following output:
\orbiteroutput{GRAPHICS/LATEX/character_table_sym_4}

\bigskip


Orbiter often uses the lexicographic order to arrange things. 
Consider the $k$-subsets of an $n$-set. 
The following example creates the $3$-subsets of a $6$-set 
in lexicographic order. 
A tree representation is produced. 
The leaves of the tree represent the subsets:
\sourcecode{all_subsets_6_3}
Here is the tree:
\orbiteroutput{GRAPHICS/LATEX/all_k_subsets_n6_k3_draw}


\bigskip


The following command illustrates how to create random $k$-subsets of a set of size $n$.
In the example, we create $20$ $5$-subsets of a $10$-element set:

\sourcecode{random_k_subsets}



\bigskip

Using the lexicographic order, 
the $k$-subsets of an $n$-element set are ranked.
The following command computes the ranks 
of a number of $3$-subsets of a $10$-element set:

\sourcecode{rank_k_subsets_test}



\bigskip

Orbiter can create the Sylvester type Hadamard matrix 
of size $2^n$ (also called the Walsh matrix). 
The following command creates the matrix of size 
$2^4 \times 2^4$ and produces a graphical representation:

\sourcecode{Walsh_matrix_4}



\bigskip

The following command creates the matrix of Dedekind numbers of order at most 10:

\sourcecode{Dedekind_10_10}

The following table is produced:
\orbiteroutput{GRAPHICS/LATEX/Dedekind_2_10_2_10}



\bigskip

The following command computes the $q$-binomial
$$
\left[
{8 \atop 3}
\right]_2
$$

\sourcecode{q_binomial_8_3_2}


\bigskip



The following command creates the elementary symmetric functions in $4$ variables.

\sourcecode{elementary_symmetric_functions_4}

The output is:
\orbiteroutput{GRAPHICS/WRAPPERS/elementary_symmetric_functions_4}

\bigskip

Orbiter can compute Domino portraits. 
To do so, we need an input file in r/g/b format 
of size $(D+1)s \times Ds,$ where $D=7$ for double-six dominos.

\sourcecode{domino_portrait}

The portrait is shown 
below
\orbiteroutput{GRAPHICS/WRAPPERS/domino_Munch}
It is possible to compare the domino portrait 
with a grayscale version of the input image.
The following command creates a grayscale image of the input file 
that was written during the previous command.

\sourcecode{domino_portrait_input}

The grayscale version of the input file is 

\orbiteroutput{GRAPHICS/WRAPPERS/domino_Munch_3_draw}





\bigskip


